% ============================================================================
% GÉNÉRÉ par scripts/exemples-guide.ts — NE PAS ÉDITER À LA MAIN.
% Régénérer : npm run exemples (chaque chiffre est vérifié contre le moteur).
% ============================================================================
\pexemples Accumulation sur le revenu de \emph{travail} (au-delà d'une exclusion),
réduction de \pc{20} sur l'$\AFNI$ — diminué, pour un couple,
de l'exemption du second revenu de travail (jusqu'à
\D{16386}).

\medskip\noindent\textbf{M3 --- personne vivant seule, 30~ans, revenu de travail de \D{15000}.}
Personne seule sans enfant (taux d'accumulation de
\pc{37.3}, maximum de \num{3812.06}) :
\begin{align*}
\text{accumulation} &= \min\bigl(\pos{\num{15000} - \num{2400}} \times \pc{37.3},\ \num{3812.06}\bigr) = \num{3812.06} \\
\text{réduction} &= \pos{\num{14885} - \num{14170.05}} \times \pc{20} = \num{142.99} \\
ACT &= \num{3812.06} - \num{142.99} = \num{3669.07}
\end{align*}
L'accumulation est \emph{plafonnée} ; la réduction commence à peine :
\D{3669.07}.

\medskip\noindent\textbf{M7 --- couple, 45 et 44~ans, revenus de travail de \D{30000} et \D{0}, sans enfant.}
Couple sans enfant, un seul revenu : le conjoint sans revenu n'apporte
aucune exemption de second revenu.
\begin{align*}
\text{accumulation} &= \min\bigl(\num{26400} \times \pc{37.3},\ \num{5943.38}\bigr) = \num{5943.38} \\
\text{réduction} &= \pos{\num{29735} - 0 - \num{21787.19}} \times \pc{20} = \num{1589.56} \\
ACT &= \num{5943.38} - \num{1589.56} = \num{4353.82}
\end{align*}
ACT : \D{4353.82}.

\medskip\noindent\textbf{M8 --- couple, 38 et 36~ans, revenus de travail de \D{60000} et \D{40000}, deux enfants : 4~ans (garde non subventionnée, \D{13000}) et 8~ans (garde non subventionnée, \D{3000}).}
Couple avec enfants : maximum de \num{3808.23} (vite atteint sur
\num{100000} de revenu de travail), mais l'exemption du second revenu
(\num{16386}, le revenu de travail du conjoint le moins payé étant net
de la cotisation RRQ supplémentaire) ne suffit pas :
\begin{align*}
\text{réduction} &= \pos{\num{86070} - \num{16386} - \num{22007.75}} \times \pc{20} = \num{9535.25}
\end{align*}
La réduction excède l'accumulation : ACT \emph{éteinte} (\D{0}).

\medskip\noindent\textbf{M4 --- personne vivant seule, 40~ans, revenu de travail de \D{50000}.}
Accumulation au maximum (\num{3812.06}), mais l'$\AFNI$ de
\num{49535} donne une réduction de \num{7072.99} : ACT éteinte
(\D{0}).
